% =======================================================================
\section{Methods}

\subsection{Data}

Data for this study come from ten professional football matches
(A-League, 2024/25 season), tracked via SkillCorner's broadcast-derived
system at 10~Hz \cite{SkillCorner2024}. Full data processing and the
topological method itself are described in a companion paper
\cite{paperA}; here, data are used only as the basis for the four
analyses below, conducted on 1{,}500 frames sampled uniformly across
the ten matches (150 per match) unless otherwise stated.

\subsection{Topological Analysis Pipeline}

Persistent homology quantifies the shape of spatial data by tracking
how topological features appear and disappear as a distance threshold
is gradually relaxed \cite{paperA}. Applied to a set of player
positions, the method builds a growing sequence of proximity networks.
At a small threshold, only closely positioned players are connected. As
the threshold increases, more connections form, and loops in the network
close temporarily before being filled. A feature's persistence, the
difference between the threshold at which it appears and the threshold
at which it disappears, measures how stable that feature is across
spatial scales. A loop with high persistence indicates a genuine
ring-like gap in player positioning that is stable over a range of
distances; a loop with low persistence is transient.

The analysis pipeline used here proceeds in three steps. Players are
first grouped by proximity at a validated threshold distance, reducing
the 22-player cloud to a smaller set of cluster centroids representing
organisational units at a specific spatial scale. A Vietoris--Rips
filtration is then computed on these centroids, producing a persistence
diagram that records the birth and death of two types of feature:
connected components ($H_0$, counting distinct spatial groups) and
loops ($H_1$, detecting ring-like gaps). $H_1$ total persistence, the
sum of loop lifespans in the diagram, provides a continuous scalar
measure of ring-like structure at that scale and moment.

The pipeline is run at two spatial scales, validated in the companion
paper \cite{paperA}: an individual scale, detecting ring-like
arrangements of players within approximately three metres of one
another, and a tactical scale, detecting larger ring-like gaps such as
a pressing encirclement, a peripheral possession orbit, or a gap
between defensive and midfield lines. The companion paper establishes
that both scales are stable across all ten matches and robust to
parameter choices. This paper treats them as pre-validated measures and
asks what they reveal about football.

\subsection{Event Annotation and Correlation Methodology}

Each match event annotated by SkillCorner, including player possessions, on-ball engagements, off-ball runs, passing options, and phases of play such as build-up, transition, and set play, was matched to the topological
time series in a window immediately surrounding it. Persistence was
averaged over the five frames before and the five frames after each
event, 0.5 seconds at the 10~Hz sampling rate used here, and the
difference between the two averages gave the persistence change
associated with that event. This narrow window captures the immediate
response to an event while limiting contamination from unrelated play. Significance of each
event type's persistence change was assessed against zero using the
Mann--Whitney $U$ test, with correction for testing multiple event
types.

To check that this pattern did not depend on the particular choice of
window width, the same comparison was repeated at three half-widths:
$\pm0.5$~s (the default), $\pm1$~s, and $\pm5$~s, for the five event
types with the clearest effect at the default width.

\subsection{Geometric Baseline Descriptors}

Four geometric descriptors already used in football analytics were
computed for every frame across all ten matches, alongside
tactical-scale $H_1$ total persistence: team length, the range of
player $x$-coordinates; team width, the range of player
$y$-coordinates; convex-hull area, the area of the smallest convex
polygon enclosing all 22 players; and Voronoi dispersion entropy, a
measure of how evenly players are spread across the pitch
\cite{shape_descriptors}. Spearman rank correlation quantified the
pairwise relationship between each descriptor and tactical-scale
persistence. Partial $R^2$, the variance in each geometric descriptor explained by persistence after accounting for the other three descriptors, quantified whether persistence added information beyond
what the geometric block already captured.

\subsection{Bilateral (Home/Away) Decomposition}

The 22-player analysis merges home and away players into a single point
cloud, which discards which team each player belongs to. To recover
team identity, the same tactical-scale pipeline, $\delta=12.0$~m with the adaptive filtration described in the companion paper \cite{paperA}, was run separately on each team's 11 players, using
the home/away labels provided with the tracking data. This gives a
separate persistence sequence for each team across the match, allowing
two questions to be asked: how often each team's own shape forms a
tactical-scale loop, and whether the two teams' sequences move together
or independently. The second question was assessed using Spearman
correlation between the two teams' persistence values at the same point
in time and at short delays of up to one second in either direction,
and by the proportion of frames in which both teams, one team, or
neither team showed a loop.

\subsection{Predictive Utility / Cross-Validated Classification}

To test whether tactical-scale persistence improves prediction rather
than simply correlating with existing measures, a classifier was
trained to detect build-up phases, defined as frames flagged as such in SkillCorner's possession-phase annotation, present in 13.7\% of frames, from frame-level features. Two feature sets were compared: the four
geometric descriptors alone, and the same four descriptors plus
tactical-scale $H_1$ total persistence. A logistic regression classifier
was trained and evaluated using ten-fold cross-validation grouped by
match, ensuring every match held out for testing was entirely absent
from training; this ensures performance reflects how well the model
generalises to a new match rather than to new frames within matches
already seen. Model performance was summarised by the area under the
ROC curve, with the difference between the two feature sets assessed
using match-level bootstrap confidence intervals and a permutation
test. The same comparison was repeated using loop count and maximum
persistence in place of total persistence, and using a second target,
chaotic-phase frames (7.0\% prevalence), to check the result was not
specific to one feature definition or one target.
